\chapter{Discussion}\label{ch:discussion}

This chapter interprets the results of Chapter~\ref{ch:results} in
relation to the two headline research questions (RQ1, RQ2) and the
construction-validity diagnostic (DQ1), then examines construct
validity, threats to validity, and limitations.

\section{RQ1: Arena Ranking Validity}\label{sec:disc-rq1}

\section{RQ2: The Migration Effect}\label{sec:disc-rq2}

\section{Construction Validity (DQ1) in Context}\label{sec:disc-dq1}

\section{Construct Validity and Threats to Validity}\label{sec:disc-validity}

\subsection{Internal Validity}\label{subsec:disc-internal-validity}

\subsection{External Validity}\label{subsec:disc-external-validity}

\subsection{Construct Validity}\label{subsec:disc-construct-validity}

\section{Limitations}\label{sec:disc-limitations}\label{sec:discussion-limitations}

% TODO: integrate with results --- the three judge- and
% statistics-facing limitations below were moved here from
% Chapter 3 (Architectural Limitations) and must be discussed
% against the observed arena outcomes; the four purely
% architectural limitations remain in Section~\ref{sec:arch-limits}.

\textbf{Judge-side benchmark contamination.}
MMLU-Pro is a public benchmark, and its questions and answer keys
plausibly appear in the pre-training corpora of contemporary LLMs;
targeted probing has demonstrated recoverable benchmark content in
modern models~\parencite{deng2024contamination}, and
contamination-limited benchmark designs exist precisely because the
problem is pervasive~\parencite{white2024livebench}.
For this thesis the threat is specific: not that the \emph{evaluated}
models saw the questions (that affects the MMLU-Pro accuracy oracle
and the arena symmetrically), but that the \emph{judge} may recognise
a held-out question and its correct answer from training.
A memorising judge approximates the accuracy oracle directly, so
judge contamination \emph{inflates} the RQ1 rank agreement for
reasons unrelated to the evaluation mechanism.
The answer-key-blind protocol does not defend against it: the
information is in the judge's weights, not in the prompt.

The generic-judge condition C3---a shared structural template in the
style of MT-Bench prompting~\parencite{zheng2023judging}, with no
leaf-specific rubric---partially disciplines this threat.
Because C1 and C3 share the same judge model, any contamination
contributes to both conditions alike, so the \emph{contrast}
$\tau_{\text{C1}} - \tau_{\text{C3}}$ is contamination-robust
evidence about the marginal value of reference-informed
conditioning: a rubric gain that survives the control cannot be
explained by memorised answer keys, which the generic judge also
possesses.
What C3 does \emph{not} do is bound the absolute level of either
condition: if contamination is strong, both conditions report
inflated agreement, and the absolute RQ1 claim inherits that
inflation.
The absolute agreement numbers are therefore interpreted as an upper
bound on mechanism-attributable validity, while the C1--C3 contrast
carries the contamination-robust part of the claim.
A direct estimate of judge memorisation---for example, verdict
consistency under paraphrased questions or a probe set published
after the judge's training cutoff---is identified as future work.

\textbf{Residual judge bias.}
The dual-call protocol mitigates positional bias but leaves two
biases architecturally unmitigated.
\textit{Verbosity bias}: LLM judges tend to prefer longer responses
independent of quality~\parencite{chen2024verbosity};
the rubric includes no explicit length-normalisation criterion.
\textit{Self-preference bias}: when the judge model coincides with an
evaluated model, verdicts may favour the judge's own response
style~\parencite{panickssery2024llm}; no judge-exclusion constraint
is enforced.
Their magnitude is bounded indirectly via the rank-stability analysis
across different judge temperatures; a direct comparison against
human preference annotations is identified as future work.

\textbf{Tie approximation and dual-call budget accounting.}
Ties are encoded as half-wins within the standard BT likelihood,
which does not include Davidson's tie-propensity
parameter~\parencite{davidson1970ties}; the approximation slightly
underestimates tie frequency.
Additionally, each logical comparison produces two API calls under the
dual-call design; budget reporting distinguishes
$B_{\text{logical}}$ (comparisons) from $B_{\text{API}} = 2 \cdot
B_{\text{logical}}$ (judge calls).
Convergence curves and efficiency comparisons are expressed in
$B_{\text{logical}}$ to remain comparable with single-call systems.

\textbf{Sparse-leaf risk at deep levels.}
A leaf with a small held-out query pool may exhaust its comparison
budget before the connectivity bootstrap achieves a fully connected
comparison graph, since each query can serve multiple pairs but the
number of pairs grows as $\binom{M}{2}$.
Such leaves will terminate in state \texttt{BUDGET\_EXHAUSTED};
their BT estimates enter the upward aggregation at full weight, which
may inflate uncertainty at the ancestor level.
The query-weighted convergence fraction $f_{\text{conv}}$ partially
absorbs this: sparse leaves contribute little weight to the global
condition.

\section{Chapter Summary}\label{sec:disc-summary}
